\documentclass[11pt]{article}
\usepackage[margin=1in]{geometry}
\usepackage{setspace}
\usepackage{amsmath,amssymb,amsthm,mathtools,bm}
\usepackage{booktabs,threeparttable,tabularx,array,multirow}
\usepackage{graphicx,float,placeins}
\usepackage{xcolor}
\usepackage{microtype}
\usepackage{enumitem}
\usepackage{natbib}
\usepackage{hyperref}
\hypersetup{colorlinks=true,linkcolor=black,citecolor=blue,urlcolor=blue}
\onehalfspacing

\newtheorem{theorem}{Theorem}
\newtheorem{proposition}{Proposition}
\newtheorem{corollary}{Corollary}
\newtheorem{lemma}{Lemma}
\theoremstyle{definition}
\newtheorem{definition}{Definition}
\newtheorem{assumption}{Assumption}
\theoremstyle{remark}
\newtheorem{remark}{Remark}

\newcommand{\Prob}{\mathbb P}
\newcommand{\E}{\mathbb E}
\newcommand{\R}{\mathbb R}
\newcommand{\ind}{\mathbf 1}
\newcommand{\cA}{\mathcal A}
\newcommand{\cB}{\mathcal B}
\newcommand{\cF}{\mathcal F}
\newcommand{\cQ}{\mathcal Q}
\newcommand{\cR}{\mathcal R}
\newcommand{\argmax}{\operatorname*{arg\,max}}
\newcommand{\argmin}{\operatorname*{arg\,min}}

\title{\textbf{When Should Investors Trust a Machine-Learned Stock Ranking?}}
\author{Enrico Bignozzi\\Swiss Finance Institute and Universit\`a della Svizzera italiana}
\date{August 2026}

\begin{document}
\maketitle

\begin{center}
\fcolorbox{black}{gray!8}{\parbox{0.91\textwidth}{\small
\textbf{Research-draft status.} The economic framework, theory specialization, empirical design, and replication code are complete. Licensed WRDS stock-level estimates are intentionally not populated in this version. Every numerical figure currently distributed with the repository is labelled as a synthetic software validation and is not evidence about U.S. equities.}}
\end{center}

\begin{abstract}
Machine-learning models rank every stock, even when predicted-return differences are too fragile to support an investment decision. We introduce a model-agnostic framework that converts a total stock ranking into a certified partial order. A relation between two stocks is retained only when their forecast-score bands are separated; otherwise the investor abstains. This construction yields a new state variable, \emph{reliable breadth}, measuring the share of cross-sectional distinctions that can be used under a prespecified decision-risk budget. We map certified relations into portfolios and show that abstention is optimal in a simple relative-value problem with estimation ambiguity and transaction costs. Statistical certification is difficult because investors observe one persistent financial history rather than independent repetitions. We therefore certify a bounded monotone envelope of pairwise ranking error with a self-normalized procedure that requires neither a mixing-rate estimate nor a long-run variance estimator. The guarantee applies to prequential forecasts generated by a predetermined finite-memory learning algorithm. The empirical design uses the preconstructed Jensen--Kelly--Pedersen stock-characteristic and return panel on WRDS, supplemented with CRSP daily cost inputs, decomposes machine-learned positions into certified and uncertified components, and evaluates how reliable breadth varies across time, stocks, models, and implementation costs. The finance contribution is the economic distinction between nominal and reliably tradable cross-sectional information; risk certification is the inferential tool used to measure it.
\end{abstract}

\noindent\textbf{Keywords:} machine learning; cross-sectional return prediction; stock ranking; abstention; partial orders; transaction costs; risk certification.

\section{Introduction}

Machine-learning models rank every stock. Investors should not trade every rank difference.

A return-prediction model maps the information available at date $t$ into forecasts
\[
\widehat\mu_{1,t},\ldots,\widehat\mu_{N_t,t}.
\]
Portfolio formation then converts these forecasts into a total order, even when two predicted returns differ by only a few basis points. The ranking has no ``I do not know'' state. A stock must be placed above or below every other stock, and a rank-based strategy can turn each mechanical distinction into capital and turnover.

This paper asks when a predicted cross-sectional ordering is reliable enough to trade. The object of interest is not forecast accuracy alone and is not a confidence interval around each individual expected return. It is the risk of the investment decision induced by a rank relation. We distinguish
\[
\text{predicted ranking}
\qquad\text{from}\qquad
\text{reliable ranking}.
\]
The first is necessarily complete. The second can be a partial order: for some pairs the evidence supports $i\succ j$, for others it supports $j\succ i$, and for the remaining pairs the investor abstains.

This distinction matters because machine learning has made high-dimensional cross-sectional return prediction a canonical empirical-asset-pricing problem. \citet{GuKellyXiu2020} compare a broad set of linear and nonlinear methods using a large stock-characteristic panel and show the importance of nonlinear interactions. The subsequent literature increasingly evaluates the investment consequences of machine-learned forecasts. \citet{Allena2026} derives ex-ante confidence intervals for risk-premium forecasts and exploits cross-sectional differences in forecast precision, while \citet{JensenKellyMalamudPedersen2026} studies portfolio learning on an implementable, net-of-cost frontier. These developments raise the bar. ``Machine learning plus uncertainty'' is not by itself a contribution. Our object is different: forecast precision is an input, while ranking reliability is the expected loss of the pairwise decisions the investor actually deploys.

We introduce a nested family of partial rankings. Let $\widehat s_{i,t}>0$ be a predetermined scale summarizing the local difficulty of using stock $i$'s forecast. At abstention threshold $q\geq0$, stock $i$ is declared superior to stock $j$ only when
\[
\widehat\mu_{i,t}-q\widehat s_{i,t}
>
\widehat\mu_{j,t}+q\widehat s_{j,t}.
\]
At $q=0$, absent ties, the relation coincides with the model's total order. As $q$ increases, small or locally noisy forecast differences are removed. The surviving relation is a strict partial order, and the family is nested.

This construction gives a new measure of the investment opportunity set. We define \emph{reliable breadth} as the fraction of economically relevant pairwise rank relations retained at the selected threshold. Nominal breadth counts every distinction generated by the model. Reliable breadth counts only distinctions that the investor has enough evidence to monetize under a prespecified error budget. The time series of reliable breadth can therefore reveal when the cross-section is easy or difficult to rank and where the reliable information is located.

The statistical problem is unusual. For each candidate threshold, the investor observes one monthly sequence of cross-sectional ranking losses. These losses inherit persistence from expected returns, volatility, liquidity, common shocks, and the forecasting algorithm itself. Treating the months as independent can understate uncertainty precisely when the investment environment is most persistent. At the same time, estimating a mixing coefficient or a long-run covariance separately for every candidate threshold is unstable and operationally unattractive.

We use the risk-certification framework developed in the companion paper \citet{Bignozzi2026SNRCPS}. The natural selected-pair error rate is not automatically monotone as abstention increases: removing easy and difficult pairs can change its denominator in either direction. We therefore certify its smallest monotone majorant. For candidate thresholds $q_1<\cdots<q_K$, let $e_t(q_k)$ be the error rate among active relations in month $t$, and define
\[
L_t(q_k)
=
\max_{\ell\geq k}e_t(q_\ell).
\]
Then $L_t(q_k)$ is bounded, nonincreasing in $k$, and dominates the economically transparent loss $e_t(q_k)$. Self-normalization uses the complete certification path to select the least-abstaining threshold whose future risk is below $\alpha$ with confidence $1-\delta$. It requires neither a block length nor a long-run variance estimate.

A practical finance implementation cannot freeze one neural network for decades. We instead certify a predetermined prequential learning rule. Each forecast is generated using a fixed-length trailing window and no future data. The complete monthly loss is therefore a measurable finite-memory function of the underlying financial process. Under the sequential limit condition stated below, the certificate controls the future decision risk of the algorithm--threshold pair, while the model can continue to update according to its predeclared rule.

We then connect abstention to portfolio choice. Each certified relation can be represented as a relative-value edge. Aggregating edge out-degree minus in-degree gives a dollar-neutral portfolio. A fixed-denominator representation yields an exact decomposition of the raw machine-learned portfolio into certified and removed edges before leverage normalization. In a simple robust relative-value problem, the optimal position is zero whenever the lower end of the predicted spread fails to cover proportional trading costs. Thus, a partial ranking is not a statistical cosmetic: it is the discrete counterpart of an economically optimal no-trade region.

The empirical analysis is organized around four questions.

\begin{enumerate}[leftmargin=2em]
\item \textbf{How much?} What share of the total machine-learned ranking is reliably distinguishable, including within the economically relevant tails?
\item \textbf{When?} How does reliable breadth vary with aggregate volatility, market liquidity, funding conditions, and common-return variation?
\item \textbf{Where?} Which stocks, characteristics, sectors, and forecast disagreements account for reliable and unreliable relations?
\item \textbf{Why does it matter?} How much gross return, turnover, cost, capacity, and net economic value comes from certified versus uncertified rank relations?
\end{enumerate}

The baseline data are the preconstructed US stock-month panel of \citet{JensenKellyPedersen2023} distributed through WRDS, using the published 153-characteristic taxonomy and its already aligned next-month excess return. CRSP daily price and volume supplement the panel for capacity-aware cost measurement. Forecasts are generated recursively by ridge, elastic net, tree, boosting, neural-network, and ensemble models. The headline portfolio evidence is net of linear and capacity-aware trading costs, and every result is reported both with endogenous exposure reduction and after equal-gross rescaling.

The paper contributes at three levels. First, it introduces reliable breadth and partial stock rankings as economic objects. Second, it provides an exact certified-versus-uncertified decomposition of rank-based positions and a no-trade interpretation. Third, it brings dependence-robust decision-risk certification to cross-sectional portfolio choice without presenting the statistical method as the finance contribution. The intended finding is not that a new algorithm forecasts returns better. It is whether machine learning identifies more apparent cross-sectional differentiation than investors can reliably and economically use.

\section{From Total Rankings to Certified Partial Orders}
\label{sec:ranking}

\subsection{Prequential forecasts and score bands}

At month $t$, let $\cA_t=\{1,\ldots,N_t\}$ be the eligible stock universe. A predetermined learning algorithm produces the next-month return forecast $\widehat\mu_{i,t}$ using information available at $t$. The algorithm may be linear, tree based, or a neural network. The primary implementation uses a fixed-length rolling training window, so the forecast sequence is genuinely out of sample and has finite memory.

Let $\widehat s_{i,t}>0$ be a reliability scale. It can be learned from prequential absolute residuals, model disagreement, or both, but it is computed without using $R_{i,t+1}$. The paper does not interpret
\[
[\widehat\mu_{i,t}-q\widehat s_{i,t},
  \widehat\mu_{i,t}+q\widehat s_{i,t}]
\]
as a stock-specific confidence interval. The bands are a nested device for deciding which relative comparisons the investor is willing to assert. Their calibration concerns the realized ranking loss of the complete decision rule.

For $q\geq0$, define
\begin{equation}
\underline\mu_{i,t}(q)
=
\widehat\mu_{i,t}-q\widehat s_{i,t},
\qquad
\overline\mu_{i,t}(q)
=
\widehat\mu_{i,t}+q\widehat s_{i,t}.
\label{eq:bands}
\end{equation}

\begin{definition}[Certified rank relation]
For distinct stocks $i,j\in\cA_t$, write
\begin{equation}
i\succ_{q,t}j
\quad\Longleftrightarrow\quad
\underline\mu_{i,t}(q)
>
\overline\mu_{j,t}(q).
\label{eq:relation}
\end{equation}
When neither $i\succ_{q,t}j$ nor $j\succ_{q,t}i$ holds, the rule abstains on the pair.
\end{definition}

Let $A_{ij,t}(q)=\ind\{i\succ_{q,t}j\}$. The next proposition gives the basic geometry.

\begin{proposition}[Nested strict partial rankings]
\label{prop:partial-order}
For every $t$ and $q\geq0$, the relation $\succ_{q,t}$ is irreflexive, asymmetric, and transitive. Moreover, if $q'\geq q$, then
\[
A_{ij,t}(q')
\leq
A_{ij,t}(q)
\]
for every pair $(i,j)$.
\end{proposition}

The result means that abstention never reverses a rank relation. It only deletes relations that are not sufficiently separated. Transitivity follows from interval separation, which is why we use stock-level bands instead of an arbitrary pairwise margin that could generate cycles.

\subsection{Reliable breadth}

Let $\omega_{ij,t}=\omega_{ji,t}\geq0$ be a pair weight measurable at $t$. The baseline uses equal weights; alternatives emphasize top-versus-bottom relations, market capitalization, or trade capacity. Define the total weight of the raw ranking by
\[
W_t
=
\sum_{i\neq j}
\omega_{ij,t}A_{ij,t}(0).
\]
Absent forecast ties, exactly one direction is active for each unordered pair with positive weight.

\begin{definition}[Reliable breadth]
The reliable breadth of the cross-section at threshold $q$ is
\begin{equation}
\cB_t(q)
=
\frac{
\sum_{i\neq j}\omega_{ij,t}A_{ij,t}(q)
}{W_t},
\label{eq:reliable-breadth}
\end{equation}
with $\cB_t(q)=0$ when $W_t=0$.
\end{definition}

Reliable breadth lies in $[0,1]$, equals one at $q=0$ under no ties, and is nonincreasing in $q$. It measures relations, not securities. A cross-section with one thousand stocks can have high nominal breadth and low reliable breadth if most forecast differences are small relative to their local scale.

The empirical analysis reports several versions. All-pair breadth measures the entire ordering. Tail breadth restricts $\omega_{ij,t}$ to pairs in which one stock lies in the raw top quantile and the other in the raw bottom quantile. Rank-distance breadth groups pairs by the distance between their raw ranks. These distinctions determine whether unreliability is confined to the middle of the cross-section or reaches the positions that drive conventional long--short strategies.

\subsection{Pairwise decision losses}

The realized error of an active relation is
\[
M_{ij,t+1}
=
\ind\{R_{i,t+1}\leq R_{j,t+1}\}.
\]
For candidate $q$, define the selected-pair error rate
\begin{equation}
e_t(q)
=
\frac{
\sum_{i\neq j}\omega_{ij,t}A_{ij,t}(q)M_{ij,t+1}
}{
\sum_{i\neq j}\omega_{ij,t}A_{ij,t}(q)
},
\label{eq:selected-error}
\end{equation}
with $e_t(q)=0$ when no relation is active. This is the natural answer to: among the comparisons the investor was willing to use, what fraction were wrong?

The denominator in \eqref{eq:selected-error} changes with $q$, so $e_t(q)$ need not be monotone. Directly applying an ordered risk-controlling selector would therefore be invalid. For an increasing finite candidate grid $q_1<\cdots<q_K$, define
\begin{equation}
L_t(q_k)
=
\max_{\ell\geq k}e_t(q_\ell).
\label{eq:monotone-envelope}
\end{equation}
Then $0\leq e_t(q_k)\leq L_t(q_k)\leq1$ and $L_t(q_k)$ is nonincreasing in $k$. We certify the stationary risk
\begin{equation}
\cR(q_k)
=
\E[L_t(q_k)].
\label{eq:ranking-risk}
\end{equation}
Controlling \eqref{eq:ranking-risk} controls the expected selected-pair error at the deployed threshold and protects against accidental nonmonotonicity at stricter thresholds.

A secondary loss avoids the random denominator:
\begin{equation}
L_t^{\mathrm{mass}}(q)
=
\frac{
\sum_{i\neq j}\omega_{ij,t}A_{ij,t}(q)M_{ij,t+1}
}{W_t}.
\label{eq:wrong-mass}
\end{equation}
It is exactly nonincreasing in $q$ and measures the weighted mass of wrong bets per raw opportunity. The main text emphasizes \eqref{eq:selected-error} because its interpretation is closer to ranking reliability; \eqref{eq:wrong-mass} is a robustness target.

\section{From Certified Relations to Positions}
\label{sec:portfolio}

\subsection{The edge representation}

Every active relation can be interpreted as a pair trade: long the predicted winner and short the predicted loser. Aggregate these edges through
\begin{equation}
\widetilde w_{i,t}(q)
=
\frac{1}{W_t}
\sum_{j\neq i}
\omega_{ij,t}
\left[
A_{ij,t}(q)-A_{ji,t}(q)
\right].
\label{eq:edge-portfolio}
\end{equation}
The portfolio is dollar neutral because the contribution of every edge sums to zero across its two endpoints. Its gross exposure declines endogenously when the investor abstains. This variable-notional portfolio measures the direct economic consequence of refusing uncertain distinctions.

For an equal-gross comparison, define
\[
w_{i,t}^{\mathrm{EG}}(q)
=
G
\frac{\widetilde w_{i,t}(q)}{
\sum_j|\widetilde w_{j,t}(q)|
}
\]
when the denominator is positive, where $G$ is the gross-leverage target. The variable-notional and equal-gross versions answer different questions. The first includes the decision to hold cash; the second isolates whether the surviving relations contain better information per unit of exposure.

The fixed denominator in \eqref{eq:edge-portfolio} gives an exact decomposition before rescaling:
\begin{equation}
\widetilde{\bm w}_t(0)
=
\widetilde{\bm w}_t(\widehat q_t)
+
\left[
\widetilde{\bm w}_t(0)
-
\widetilde{\bm w}_t(\widehat q_t)
\right].
\label{eq:exact-decomposition}
\end{equation}
The first component is certified and the second is the portfolio generated by removed relations. Equation \eqref{eq:exact-decomposition} permits a clean attribution of returns and changes in positions. Trading-cost attribution additionally accounts for the fact that costs are nonlinear in changes of the combined portfolio; we therefore report stand-alone costs, marginal costs, and a Shapley decomposition as robustness.

The paper also reports conventional top-minus-bottom decile portfolios. Those strategies are familiar benchmarks, but they do not give an exact relation-level decomposition. The edge portfolio is the primary economic experiment because it maps the certified object directly into capital.

\subsection{A no-trade interpretation}

Consider a single predicted relative-value opportunity. Let
\[
d_{ij,t}
=
\widehat\mu_{i,t}-\widehat\mu_{j,t}
\geq0,
\qquad
s_{ij,t}
=
\widehat s_{i,t}+\widehat s_{j,t}.
\]
The investor considers a nonnegative position $x$ that is long $i$ and short $j$. The unknown expected spread belongs to the ambiguity interval
\[
[d_{ij,t}-q s_{ij,t},
  d_{ij,t}+q s_{ij,t}].
\]
Let $c_{ij,t}\geq0$ be the proportional round-trip cost per unit of the pair position, and let $\gamma\nu_{ij,t}>0$ be the quadratic risk penalty.

\begin{proposition}[Certified ranking and the no-trade region]
\label{prop:no-trade}
The robust position
\begin{equation}
x_{ij,t}^{\star}(q)
=
\argmax_{x\geq0}
\min_{m\in[d_{ij,t}-q s_{ij,t},\,d_{ij,t}+q s_{ij,t}]}
\left\{xm-c_{ij,t}x-\frac{\gamma\nu_{ij,t}}{2}x^2\right\}
\label{eq:robust-pair-choice}
\end{equation}
is
\begin{equation}
x_{ij,t}^{\star}(q)
=
\frac{
[d_{ij,t}-q s_{ij,t}-c_{ij,t}]_+
}{\gamma\nu_{ij,t}}.
\label{eq:no-trade-solution}
\end{equation}
Consequently, the investor does not trade the predicted ranking relation whenever
\begin{equation}
d_{ij,t}
\leq
q s_{ij,t}+c_{ij,t}.
\label{eq:no-trade-region}
\end{equation}
\end{proposition}

Without costs, condition \eqref{eq:no-trade-region} contains the score-band overlap rule. Trading costs expand the no-trade region. The proposition is intentionally simple: it does not claim that every institutional portfolio solves independent pair problems. Its role is to show that abstention has a standard economic interpretation. Weak forecast differences should not automatically become positions when their worst-case spread is nonpositive net of implementation costs.

\subsection{Costs and capacity}

Let $\Delta w_{i,t}=w_{i,t}-w_{i,t-1}^{+}$ denote the trade after accounting for passive weight drift. We report a transparent linear benchmark
\[
\mathrm{TC}_{t}^{\mathrm{lin}}
=
c\sum_i|\Delta w_{i,t}|
\]
and a capacity-aware specification
\begin{equation}
\mathrm{TC}_{t}^{\mathrm{cap}}
=
\sum_i
\left[
\frac{\mathrm{spread}_{i,t}}{2}|\Delta w_{i,t}|
+
\lambda\sigma_{i,t}^{d}
\sqrt{
\frac{\mathrm{AUM}|\Delta w_{i,t}|}
{\mathrm{ADV}_{i,t}}
}
|\Delta w_{i,t}|
\right].
\label{eq:capacity-cost}
\end{equation}
The inputs use information available at portfolio formation. We trace results over AUM and risk rather than reporting one cost assumption. This follows the economic principle in \citet{JensenKellyMalamudPedersen2026}: a strategy should be judged on the attainable net-return frontier, not on gross alpha alone.

\section{Risk Certification from One Financial History}
\label{sec:certification}

\subsection{Candidate library, certification, and deployment}

The baseline candidate library is deterministic and fixed before the empirical analysis. We use a dense increasing grid $0=q_1<\cdots<q_K$ for the dimensionless score-band multiplier. This choice eliminates dependence between library construction and the certification trajectory. It also preserves sample length: no proposal block is needed. Because ordered selection does not pay a multiplicity factor in $K$, grid refinement is statistically inexpensive.

As a robustness check, a proposal block can use the forecast-separation margins
\[
\frac{
|\widehat\mu_{i,t}-\widehat\mu_{j,t}|
}{
\widehat s_{i,t}+\widehat s_{j,t}
}
\]
to target particular reliable-breadth levels. For a rolling learner with memory $p$, this random library is separated from certification by a guard gap exceeding $p$; otherwise the proposal variables and certification losses overlap through model training. The deterministic grid is therefore the primary design. The certification block contains $m$ monthly cross-sections and produces the bounded monotone losses in \eqref{eq:monotone-envelope}.

For candidate $q_k$, define
\begin{equation}
\widehat\cR_m(q_k)
=
\frac1m\sum_{t=1}^{m}L_t(q_k)
\label{eq:empirical-risk}
\end{equation}
and the recursive self-normalizer
\begin{equation}
\widehat V_m(q_k)
=
\frac{1}{m^2}
\sum_{\ell=1}^{m}
\left[
\sum_{t=1}^{\ell}
\left[L_t(q_k)-\widehat\cR_m(q_k)\right]
\right]^2.
\label{eq:self-normalizer-finance}
\end{equation}
Let $B$ be standard Brownian motion and let $c_{1-\delta}$ be the $(1-\delta)$ quantile of
\begin{equation}
\mathsf T
=
\frac{B(1)}{
\left\{
\int_0^1[B(r)-rB(1)]^2\,dr
\right\}^{1/2}}.
\label{eq:brownian-pivot-finance}
\end{equation}
The self-normalized upper risk bound is
\begin{equation}
U_{m,k}^{\mathrm{SN}}
=
\min\left\{
1,
\widehat\cR_m(q_k)
+
c_{1-\delta}
\left(
\frac{\widehat V_m(q_k)}{m}
\right)^{1/2}
\right\}.
\label{eq:sn-ucb-finance}
\end{equation}
The selected threshold is the smallest candidate for which the bound at that candidate and every stricter candidate is below $\alpha$:
\begin{equation}
\widehat k
=
\min\left\{
k:
U_{m,\ell}^{\mathrm{SN}}\leq\alpha
\text{ for every }\ell\geq k
\right\}.
\label{eq:finance-selector}
\end{equation}
If no finite candidate passes, the rule abstains from every comparison.

The critical value in \eqref{eq:sn-ucb-finance} is $c_{1-\delta}$, not $c_{1-\delta/K}$. Ordered monotonicity means that an unsafe selected rule requires a failure at one boundary candidate, so the nominal confidence level does not deteriorate with the resolution of the candidate grid. This is the risk-controlling selection principle of \citet{BatesEtAl2021}; the dependence-robust self-normalized bound is specialized from \citet{Bignozzi2026SNRCPS}.

\subsection{Certifying a dynamic learning rule}

Let $\bm Z_t$ contain the complete monthly cross-section, returns, characteristics, and predetermined aggregate inputs. The primary forecasting algorithm uses a fixed rolling window of length $p$ and a fixed refitting protocol. Thus, for every candidate $q$, the ranking loss can be written as
\begin{equation}
L_t(q)
=
h_q(\bm Z_{t-p},\ldots,\bm Z_{t+1}),
\label{eq:finite-memory-loss}
\end{equation}
where the lead appears only because $L_t$ is realized after next-month returns become available. The algorithm is dynamic, but the map $h_q$ is fixed before proposal. Under a stationary underlying process, \eqref{eq:finite-memory-loss} defines a stationary finite-memory loss process. This is the object being certified.

\begin{assumption}[Sequential limit condition]
\label{ass:finance-fclt}
For $q$ in a compact candidate interval, the centered partial-sum process
\[
\left\{
\frac1{\sqrt m}
\sum_{t=1}^{\lfloor mr\rfloor}
[L_t(q)-\cR(q)]
:(q,r)\in\cQ\times[0,1]
\right\}
\]
converges weakly to a continuous Gaussian process whose time section at each $q$ has distribution $\sigma(q)B(\cdot)$, with $\sigma(q)$ continuous and bounded away from zero. The associated Brownian-bridge denominator is uniformly nondegenerate.
\end{assumption}

This assumption is stated directly because no single mixing exponent is sufficient for every high-dimensional learning rule. A fixed finite library only requires a multivariate functional central limit theorem. For finite-memory forecasts under geometrically absolutely regular data, standard empirical-process conditions give one primitive route.

\begin{theorem}[Certified partial ranking]
\label{thm:certified-ranking}
Suppose Assumption \ref{ass:finance-fclt} holds and the pair weights and forecasting protocol are fixed before certification outcomes are observed. The candidate grid is either deterministic or generated on a proposal block separated from certification by a guard gap $g$. Let $\rho_g=0$ for the deterministic grid and let $\rho_g$ denote the proposal--certification coupling remainder for a random grid. Let $\widehat q=q_{\widehat k}$ be selected by \eqref{eq:finance-selector}. Then
\begin{equation}
\Prob\left\{
\cR(\widehat q)>\alpha
\right\}
\leq
\delta
+
\varepsilon_m(\delta,\cQ)
+
\rho_g,
\label{eq:finance-risk-bound}
\end{equation}
where $\varepsilon_m(\delta,\cQ)\to0$. If the loss at date $t$ has memory $p$ and the underlying process is absolutely regular, one may take $\rho_g\leq\beta_{\bm Z}(g-p-1)$ whenever $g>p+1$. Moreover,
\begin{equation}
\Prob\left\{
\E[e_0(\widehat q)]>\alpha
\right\}
\leq
\delta
+
\varepsilon_m(\delta,\cQ)
+
\rho_g.
\label{eq:pairwise-error-bound}
\end{equation}
The bound contains no factor depending on $K$.
\end{theorem}

Equation \eqref{eq:pairwise-error-bound} follows because $e_t(q)\leq L_t(q)$ pointwise. The guarantee is training conditional in the sense relevant for a portfolio manager: with high probability over the historical trajectory used to choose $\widehat q$, the deployed algorithm--threshold pair has stationary decision risk below $\alpha$. It is not a conditional-coverage statement for every stock or market state.

\subsection{Repeated deployment}

The baseline recertifies annually and freezes $\widehat q$ for the next twelve deployment months. Forecasts continue to update according to the predetermined rolling algorithm. If epoch $j$ receives confidence budget $\delta_j$ and $\sum_j\delta_j\leq\delta$, a union bound gives joint control over all scheduled deployments without requiring independence between epochs. An event-triggered version can monitor realized loss or reliable breadth at deterministic checkpoints and activate a new, separately certified threshold. The trigger decides when to seek a replacement; it does not substitute for certification and does not create an arbitrary-stopping theorem.

Finite-sample exactness without a dependence envelope is impossible even for simple persistent Markov losses, as shown in the companion paper. We therefore report a blocked finite-sample benchmark when a conservative mixing envelope is imposed, but use self-normalization as the primary dependence-blind procedure. The distinction is explicit: the blocked result is finite sample and assumption heavy; the self-normalized result is tuning free and process-specific asymptotic.

\section{Data and Empirical Design}
\label{sec:data-design}

\subsection{Stock-level panel}

The baseline data are the monthly stock-level Global Factor Data of \citet{JensenKellyPedersen2023}, distributed on WRDS as \texttt{contrib.global\_factor}. This is a preconstructed panel containing security identifiers, firm characteristics, and future stock returns in a common time-stamped format. The table currently contains more than four hundred characteristic fields. We use the 153 characteristics that enter the published JKP factor taxonomy, with their names pinned to an exact release of the authors' factor-details file. This choice gives a standard, economically interpretable predictor set without reconstructing hundreds of signals or merging a moving public characteristic release to returns.

The baseline sample contains US stocks and applies the four JKP screens
\begin{equation}
\texttt{common}=1,
\qquad
\texttt{exch\_main}=1,
\qquad
\texttt{primary\_sec}=1,
\qquad
\texttt{obs\_main}=1.
\label{eq:jkp-screens}
\end{equation}
These restrictions retain common equity on a country's main exchanges, select the primary security when a firm has multiple securities, and avoid duplicate CRSP and Compustat observations. The prediction target at formation month $t$ is \texttt{ret\_exc\_lead1m}, the next-month excess return supplied by the same panel. The baseline extraction runs from January 1963 through December 2025. The precise ending month is verified against the live WRDS table and recorded in the extraction manifest.

Descriptive estimates of ranking reliability use the full screened US cross-section. Baseline investable portfolios retain the JKP mega, large, and small size groups and exclude micro and nano stocks. We restore the complete cross-section in robustness tests and report all results separately by size group. This division prevents the headline economic value from being driven mechanically by securities whose nominal forecast spreads are difficult to implement.

The JKP panel also contains market equity and several liquidity or spread characteristics. These variables support mechanism tests and reduced-form cost robustness. The primary capacity analysis supplements the monthly panel with CRSP daily price and volume data, from which we construct trailing average dollar volume and daily volatility. A ready-made predictor-and-return panel does not eliminate the need to measure implementation costs separately. The replication archive therefore records the JKP release, live WRDS schema, selected columns, standard screens, extraction dates, annual partition counts, CRSP cost-data vintage, and file checksums before any production estimate is generated.

As an external standardized benchmark, we additionally support the JKP Common Task Framework tables \texttt{contrib\_global\_factor.ctff\_chars}, \texttt{ctff\_features}, and \texttt{ctff\_daily\_ret}. They provide a compact model-input interface and a common evaluation environment. They are not mixed into the headline sample: \texttt{contrib.global\_factor} remains the baseline because it exposes the historical stock-month panel and the published characteristic taxonomy directly.

\subsection{Prequential forecasts}

For each month, every model is estimated using only a fixed-length trailing sample ending before the forecast date. The baseline training window is ten years; fifteen- and twenty-year windows are robustness checks. Firm characteristics are cross-sectionally rank transformed to $[-1,1]$ each month and missing values are set to the cross-sectional median after the transform. Industry indicators and a small set of predetermined aggregate variables enter only where specified.

The model grid follows the empirical asset-pricing literature rather than introducing a bespoke forecasting architecture:
\begin{enumerate}[leftmargin=2em]
\item ridge regression;
\item elastic net;
\item extremely randomized trees;
\item gradient-boosted trees;
\item multilayer perceptrons;
\item an equal-weight forecast ensemble.
\end{enumerate}
Baseline hyperparameters are fixed ex ante in the repository configuration. Nested chronological validation within the rolling training window is a robustness exercise rather than the source of the headline specification. The final deployment month never affects either choice.

The reliability scale $\widehat s_{i,t}$ is fitted from prequential absolute forecast residuals available before $t$. The baseline combines a residual-scale model with cross-model forecast dispersion. Alternatives use each component separately, a constant scale, realized-volatility scaling, and Bayesian forecast standard errors where available. This decomposition is important for comparison with \citet{Allena2026}: our certificate concerns realized ranking decisions, not whether a stock-specific posterior interval contains an unknown expected return.

\subsection{Certification schedule}

The prequential prediction panel is generated once and then treated as the input to the certification layer. The baseline annual schedule uses a deterministic grid $q\in\{0,0.1,\ldots,3.0\}$, one hundred twenty trailing certification months, and twelve deployment months. At the end of each deployment, the threshold is recertified. All windows move forward chronologically. The first deployable date is determined by the rolling forecast-training window plus the certification history; it is never backfilled.

The main risk budgets are $\alpha\in\{0.25,0.30,0.35,0.40\}$ at confidence $1-\delta=0.90$. We report the complete risk--breadth frontier rather than selecting $\alpha$ from deployment returns. The dense grid is fixed independently of returns and always includes a fallback that abstains from all relations. As a robustness check, a sixty-month proposal block targets specified breadth levels and is followed by a guard exceeding the complete forecasting-and-scale algorithm's effective memory before certification begins. With the baseline ten-year forecast window, five-year residual-scale history, and annual refits, the conservative implementation records a 192-month memory bound and uses a 194-month guard. This deliberately inefficient robustness design is not needed for the deterministic main grid.

\subsection{Portfolio experiments}

The core experiments compare:
\begin{enumerate}[leftmargin=2em]
\item the raw relation-graph portfolio at $q=0$ and the otherwise identical certified graph portfolio at $q=\widehat q$;
\item the conventional raw top-minus-bottom decile strategy as a separate benchmark;
\item raw fixed-denominator edge positions;
\item certified variable-notional edge positions;
\item certified edge positions rescaled to the raw gross exposure;
\item the exact removed-edge component;
\item forecast-confidence filters and simple predicted-spread cutoffs;
\item cost-aware portfolio optimization using raw and certified signals.
\end{enumerate}

Pair relations are equal weighted in the baseline and value weighted in robustness tests. The certification exercise is reported both over all stock pairs and over the economically central raw top-versus-bottom tail pairs. Portfolio turnover is computed after passive return drift over the union of consecutive stock universes, so a security that delists or fails the next formation filter is explicitly liquidated rather than silently disappearing. We impose position caps in the optimization experiments, exclude microcaps from the headline portfolios, and report factor exposures. Net returns are evaluated at linear one-way costs of 10, 25, and 50 basis points and under \eqref{eq:capacity-cost} for AUM levels from ten million to one billion dollars.

\subsection{Inference and anti-overfitting discipline}

Time-series means and differences use heteroskedasticity-and-autocorrelation robust standard errors. Strategy comparisons additionally use paired block bootstraps over months and report performance by decade. Cross-sectional mechanism regressions cluster by month and stock where appropriate. The certification confidence statement is not replaced by these conventional standard errors; the two answer different questions.

The headline hypotheses, model grid, cost grid, sample filters, and primary exhibits are fixed in the repository configuration before WRDS results are attached. Any post-results specification is labelled exploratory. This separation is essential because a paper about reliability cannot itself rely on an undisclosed specification search.

\section{Pre-Specified Empirical Findings and Exhibit Map}
\label{sec:empirical-map}

This section states the empirical tests and the claims each test could support. It deliberately contains no numerical WRDS stock-level estimates in the current research draft.

\subsection{How much of the ranking is reliable?}

The first table reports average, median, and tail quantiles of $\cB_t(\widehat q)$ by model and risk budget. The first figure plots the probability that a relation survives against raw forecast spread, normalized spread, and rank distance. Separate panels cover all pairs, top-versus-bottom decile pairs, and adjacent ranks.

The economically relevant null is not $\cB_t=1$, which is unlikely by construction at strict risk budgets. We test whether reliable breadth remains materially below nominal breadth even among extreme predicted-return portfolios. A finding confined to adjacent middle ranks would be less important than a finding that removes a substantial fraction of top--bottom relations.

\subsection{When does reliable breadth contract?}

The second figure plots reliable breadth through time and overlays aggregate volatility, market illiquidity, recessions, and funding-stress measures. Predictive regressions take the form
\[
\cB_{t+1}
=
a+b^{\prime}\bm H_t+u_{t+1},
\]
where $\bm H_t$ is predetermined. The aim is descriptive state dependence, not causal identification. We also test whether contractions in reliable breadth forecast weaker net performance of raw rank strategies and a larger certified-versus-raw return gap.

\subsection{Where is the reliable information?}

We estimate relation-level survival probabilities as functions of size, liquidity, idiosyncratic volatility, characteristic extremeness, sector, forecast disagreement, and model class. To avoid treating millions of pairs as independent observations, inference aggregates relation outcomes by month and uses stock-level and month-level clustering in supporting regressions.

A central test conditions jointly on predicted spread and scale. If reliability is only a relabelling of forecast magnitude, the scale and other state variables should add little after conditioning on the spread. If ranking reliability is distinct, economically similar predicted differences should have different survival and realized-error rates across stocks and states.

\subsection{What do uncertified trades contribute?}

Equation \eqref{eq:exact-decomposition} separates the raw edge portfolio into certified and removed components. The third table attributes gross returns, average positions, turnover, linear costs, capacity-aware costs, and net returns. The strongest economically meaningful pattern would be that removed relations account for a disproportionately large share of turnover and costs relative to persistent gross return. The opposite result is possible and would reject the proposed economic mechanism.

Because transaction costs are nonlinear, we report three attributions: stand-alone costs of each component, marginal cost from adding the component to the other, and a two-component Shapley allocation. No single attribution is presented as an accounting identity.

\subsection{Does abstention improve the implementable frontier?}

The fourth figure traces expected net return against volatility, reliable breadth, and turnover. Variable-notional certified portfolios can improve utility by reducing exposure; equal-gross certified portfolios test whether the retained relations have greater information quality. We report Sharpe ratios, certainty equivalents, maximum drawdowns, and capacity at multiple AUM levels. Economic significance requires improvement over simple spread thresholds and over forecast-confidence filtering, not merely over an unfiltered neural network.

\subsection{Robustness and falsification}

The robustness design includes:
\begin{enumerate}[leftmargin=2em]
\item alternative characteristic releases and sample endpoints;
\item NYSE breakpoints, microcap exclusions, and value weighting;
\item binary error, normalized pairwise regret, and wrong-bet-mass losses;
\item constant, residual, disagreement, volatility, and confidence-interval scales;
\item fixed rolling and expanding forecast windows;
\item scheduled and event-triggered recertification;
\item blocked finite-sample certification under a conservative dependence envelope;
\item placebo thresholds selected after permuting forecast spreads within month;
\item negative-control rankings based on stale or random characteristics.
\end{enumerate}

The paper is successful only if reliable breadth is stable across reasonable definitions and has incremental economic content. A method that selects a sparse ranking but does not alter realized error, costs, or net investment outcomes would not support the finance claim.

\section{Conclusion}

A machine-learning model has no difficulty ranking one thousand stocks. An investor's problem is deciding which of those distinctions deserve capital.

This paper replaces the mechanically complete ranking with a certified partial order. Score-band separation creates a nested family of relation graphs; self-normalized calibration selects the least-abstaining graph whose future pairwise decision risk is controlled from one dependent financial history. The resulting reliable breadth measures the usable fraction of the cross-section. The edge representation maps that object directly into positions and separates certified from removed trades, while a robust pair-choice problem shows why uncertainty and trading costs generate a no-trade region.

The framework changes the economic question from ``Which model ranks stocks best?'' to ``How much of a model's ranking is reliably and implementably different?'' The empirical design is built to answer how much, when, where, and why. It evaluates whether uncertified relations are predominantly weak alpha, expensive turnover, or valuable information that the certificate removes. Each outcome is informative. The contribution is not guaranteed outperformance; it is a disciplined measurement of the gap between apparent and tradable machine-learned differentiation.


\appendix
\section{Proofs}

\subsection{Proof of Proposition \ref{prop:partial-order}}

\begin{proof}
Irreflexivity follows because $\underline\mu_{i,t}(q)\leq\overline\mu_{i,t}(q)$. If $i\succ_{q,t}j$, then
\[
\underline\mu_{i,t}(q)
>
\overline\mu_{j,t}(q)
\geq
\underline\mu_{j,t}(q),
\]
so $j\succ_{q,t}i$ is impossible. For transitivity, suppose $i\succ_{q,t}j$ and $j\succ_{q,t}k$. Then
\[
\underline\mu_{i,t}(q)
>
\overline\mu_{j,t}(q)
\geq
\underline\mu_{j,t}(q)
>
\overline\mu_{k,t}(q),
\]
which implies $i\succ_{q,t}k$. Finally, for $q'\geq q$,
\[
\underline\mu_{i,t}(q')
\leq
\underline\mu_{i,t}(q),
\qquad
\overline\mu_{j,t}(q')
\geq
\overline\mu_{j,t}(q).
\]
Therefore, every relation active at $q'$ is also active at $q$.
\end{proof}

\subsection{Proof of Proposition \ref{prop:no-trade}}

\begin{proof}
For $x\geq0$, the worst expected spread in \eqref{eq:robust-pair-choice} is $d_{ij,t}-q s_{ij,t}$. The robust objective is therefore
\[
x[d_{ij,t}-q s_{ij,t}-c_{ij,t}]
-
\frac{\gamma\nu_{ij,t}}{2}x^2.
\]
It is strictly concave. Its unconstrained maximizer is
\[
\frac{d_{ij,t}-q s_{ij,t}-c_{ij,t}}
{\gamma\nu_{ij,t}}.
\]
Projection onto $[0,\infty)$ gives \eqref{eq:no-trade-solution}, and the position is zero exactly under \eqref{eq:no-trade-region}.
\end{proof}

\subsection{Proof of Theorem \ref{thm:certified-ranking}}

\begin{proof}
For every candidate $q_k$, the loss in \eqref{eq:monotone-envelope} is bounded in $[0,1]$ and is nonincreasing in $k$. Assumption \ref{ass:finance-fclt} and the continuous-mapping argument for the recursive self-normalizer imply the uniform one-sided calibration statement
\[
\sup_{q\in\cQ}
\left[
\Prob\left\{
\cR(q)>U_m^{\mathrm{SN}}(q;\delta)
\right\}
-
\delta
\right]_+
\leq
\varepsilon_m(\delta,\cQ),
\]
where $\varepsilon_m(\delta,\cQ)\to0$. A Berbee coupling across the proposal--certification guard gap adds $\beta(g)$ and permits the proposal-generated random grid to be treated as independent of the certification trajectory.

Conditional on the proposal grid, let $k^\dagger$ be the largest unsafe candidate, so $\cR(q_{k^\dagger})>\alpha$ and all stricter candidates are safe. The selector can deploy an unsafe rule only if the upper bound at $k^\dagger$ fails to cover its risk. Hence the probability of unsafe selection is at most the pointwise calibration failure at that boundary candidate. This gives \eqref{eq:finance-risk-bound} without a factor depending on $K$. Since $e_t(q)\leq L_t(q)$ for every $t$ and $q$, $\E[e_0(q)]\leq\cR(q)$, which gives \eqref{eq:pairwise-error-bound}.
\end{proof}

\section{Exact Pair Counting without Materializing the Relation Graph}

A monthly cross-section with $N$ stocks contains $N(N-1)/2$ raw relations. Materializing every pair is unnecessary. For a fixed threshold $q$, define
\[
a_i
=
\widehat\mu_{i,t}-q\widehat s_{i,t},
\qquad
b_j
=
\widehat\mu_{j,t}+q\widehat s_{j,t}.
\]
Stock $i$ dominates stock $j$ if $b_j<a_i$. Sorting stocks by $a_i$ and $b_j$ gives the active-pair count in $O(N\log N)$ operations. To count errors, compress realized returns into ranks and maintain a Fenwick tree over the returns of stocks whose upper endpoint has entered the active set. For each $i$, the tree returns the weight of active $j$ with $R_{j,t+1}\geq R_{i,t+1}$. Repeating this computation over $K$ candidates costs $O(KN\log N)$ and uses $O(N)$ memory.

The implementation in \texttt{src/rankcert/partial\_order.py} supports product pair weights and is checked against brute force in the unit tests. Reliable breadth and graph degree require only one-dimensional sorting and cumulative sums.

\section{Prequential Learning Protocol}

Let $p$ be the forecast-training-window length. For each month $t$:
\begin{enumerate}[leftmargin=2em]
\item form characteristics using information available at $t$;
\item fit or update the model using observations with formation dates in $[t-p,t-1]$ and realized targets available by $t$;
\item produce $\widehat\mu_{i,t}$ for the eligible cross-section;
\item produce $\widehat s_{i,t}$ from residuals and model disagreement available by $t$;
\item store the prediction, scale, universe, and version metadata;
\item reveal $R_{i,t+1}$ only after portfolio formation.
\end{enumerate}
The entire forecast panel is generated before the certification analysis selects thresholds. This prevents the certificate from feeding back into forecast-model tuning.

\section{Synthetic Software Validation}

The repository includes a persistent latent-state stock panel solely to validate the end-to-end code path. It is not calibrated to U.S. equities and is not used in any empirical claim. The validation verifies that the optional proposal-grid branch and the primary fixed-grid branch both execute, the selected-pair error is replaced by its monotone envelope, the self-normalized bound selects a nontrivial threshold, and deployment outputs include pair error, reliable breadth, turnover, costs, and net returns.

\IfFileExists{figures/concept_partial_order.pdf}{
\begin{figure}[H]
\centering
\includegraphics[width=0.78\textwidth]{figures/concept_partial_order.pdf}
\caption{Conceptual score-band construction. Overlapping score bands imply abstention. This is a schematic, not an empirical estimate.}
\end{figure}}{}

\IfFileExists{figures/certification_frontier.pdf}{
\begin{figure}[H]
\centering
\includegraphics[width=0.72\textwidth]{figures/certification_frontier.pdf}
\caption{Synthetic certification frontier used to validate the software. The plot is not evidence about stock-ranking reliability.}
\end{figure}}{}

\section{Planned Main Tables and Figures}

\begin{table}[H]
\centering
\caption{Planned headline exhibits}
\begin{tabularx}{\textwidth}{@{}p{0.12\textwidth}X@{}}
\toprule
Exhibit & Content \\
\midrule
Figure 1 & Probability a rank relation survives against predicted spread, normalized spread, and rank distance. \\
Table 1 & Reliable breadth by model, risk budget, all-pair universe, and tail-pair universe. \\
Figure 2 & Time series of reliable breadth with aggregate uncertainty and liquidity states. \\
Table 2 & Relation-level anatomy by size, liquidity, idiosyncratic volatility, sector, and model disagreement. \\
Table 3 & Certified versus removed-edge attribution of gross return, turnover, costs, and net return. \\
Figure 3 & Implementable frontier: expected net return against risk and reliable breadth at multiple AUM levels. \\
Table 4 & Comparison with forecast-confidence filtering and simple predicted-spread thresholds. \\
\bottomrule
\end{tabularx}
\end{table}


\bibliographystyle{plainnat}
\bibliography{references}
\end{document}
